\begin{definition}
For a sample \(\omega\), parameter \(\varepsilon\), and finite set \(I\), the
pre-terminal recorded projected coordinates are the tagged coordinates recorded
before the terminal Bernoulli scan.  They are the elements of
\(\noderef{TwoBiteTerminalCoordinateUniverse}(m)\) satisfying one of the two
pre-terminal recording conditions below.

First, the coordinate may be an incident coordinate revealed for a
staged-revealed vertex.  A red coordinate \(\operatorname{red}(q)\) is recorded
in this way if, for some red base vertices \(r,s\), the coordinate \(q\) is
\((r,s)\) or \((s,r)\), and either
\[
\noderef{TwoBiteStageRevealedVertex}
(\omega,\varepsilon,I,\operatorname{red}(r))
\quad\text{and}\quad
\noderef{TwoBiteProjectionContains}
(\omega,I,\operatorname{red}(s)),
\]
or the same assertion holds with \(r\) and \(s\) interchanged.  The blue branch
is identical with red replaced by blue.  This branch records the neighbor and
non-neighbor queries made for \(F_0\), \(F_1\), and \(F_2\); in particular it
contains the same-color coordinates queried while evaluating \(F_2\), because
the \(F_2\) count only asks for adjacencies to \(F_1\)-vertices, and
\(F_1\subseteq F\).

Second, after \(F=F_0\cup F_1\cup F_2\) is known, the coordinate may be
recorded by the paper's instruction to reveal all pairs in \(X_x(I)\) for
\(x\in F\).  Thus a red coordinate \(\operatorname{red}(q)\) is recorded in
this way if there is a base vertex \(x\) with
\[
\noderef{TwoBiteStageRevealedVertex}(\omega,\varepsilon,I,x)
\quad\text{and}\quad
q\in\noderef{TwoBiteBasePairs}
\bigl(\pi_R^\omega(\noderef{TwoBiteX}(\omega,I,x))\bigr).
\]
For a blue coordinate \(\operatorname{blue}(q)\), the same condition holds
with \(\pi_B^\omega\) in place of \(\pi_R^\omega\).
\end{definition}
